% Generated draft: requires mathematical and source review.
\documentclass[11pt]{article}
\usepackage[margin=25mm]{geometry}
\usepackage{fontspec}
\setmainfont{Latin Modern Roman}
\usepackage{amsmath,amssymb,mathtools,unicode-math}
\setmathfont{Latin Modern Math}
\usepackage{enumitem,needspace,xurl,hyperref}
\setlength{\parindent}{0pt}
\setlength{\parskip}{5pt}
\setlength{\emergencystretch}{3em}
\newcommand{\N}{\mathbb{N}}
\newcommand{\Z}{\mathbb{Z}}
\newcommand{\Q}{\mathbb{Q}}
\newcommand{\R}{\mathbb{R}}
\newcommand{\C}{\mathbb{C}}
\newcommand{\F}{\mathbb{F}}
\newcommand{\A}{\mathbb{A}}
\newcommand{\PP}{\mathbb{P}}
\newcommand{\E}{\mathbb{E}}
\newcommand{\eps}{\varepsilon}
\newcommand{\dd}{\,\mathrm{d}}
\begin{document}
\section*{UnsolvedMath ID: 20003266}
\textbf{Source problem code:} catalogue entry
\subsection{Definitions and mathematical statement}
This extract is a sentence fragment referring to Conjecture 3, not a separately stated conjecture. It says that ``it is not even known whether an Anosov automorphism is topologically transitive in general'' and ends with the incomplete phrase ``the conjugation map in''. The supplied verification confirms the printed word ``automorphism''; it must not be silently changed.

For the algebraic toral interpretation, let $d\ge1$, $A\in\mathrm{GL}(d,\Z)$, $\mathbb T^d=\R^d/\Z^d$, and $T_A(x)=Ax\pmod{\Z^d}$. Topological transitivity means that for every pair of nonempty open sets $U,V$, some $n\ge0$ satisfies $T_A^n(U)\cap V\ne\varnothing$. Topological mixing requires this for all sufficiently large $n$. Haar mixing means
\[
 \int_{\mathbb T^d} f(T_A^n x)g(x)\,\dd m(x)
 \longrightarrow \left(\int f\,\dd m\right)\left(\int g\,\dd m\right)
\]
for all $f,g\in L^2(m)$, where $m$ is normalized Haar measure. The supplied aid reports that, for toral automorphisms, these three properties are equivalent to $A$ having no root-of-unity eigenvalue. In particular, a hyperbolic toral automorphism, meaning $|\lambda|\ne1$ for every eigenvalue, is mixing; its transitivity is not an open question.

The aid's finite-Fourier formulation is as follows. For finite $K,L\subset\Z^d$ and no root-of-unity eigenvalue, the set
\[
 \{n\ge0:(A^{\mathsf T})^n k=-\ell
       \text{ for some }k\in K,\ \ell\in L,\ (k,\ell)\ne(0,0)\}
\]
is finite. Thus, for trigonometric polynomials $p,q$ with supports in $K,L$, their above correlation equals the product of their means exactly for all sufficiently large $n$. The zero-frequency pair is excluded from the collision set because it supplies the constant term. The aid also reports that an affine hyperbolic toral map $x\mapsto Ax+b$ is translation-conjugate to $T_A$ and inherits mixing.

A different, unresolved reading concerns arbitrary Anosov diffeomorphisms of closed connected manifolds. A smooth diffeomorphism $f:M\to M$ is Anosov if there is a continuous $Df$-invariant splitting $TM=E^s\oplus E^u$ and constants $C>0$, $0<\lambda<1$ such that for every $n\ge0$,
\[
 \|Df^n v\|\le C\lambda^n\|v\|\quad(v\in E^s),\qquad
 \|Df^{-n}w\|\le C\lambda^n\|w\|\quad(w\in E^u).
\]
The plausible contextual question is whether every such $f$ is topologically transitive. The aid reports this general assertion as still open in the primary literature it checked, but the extract does not itself supply a complete statement with regularity, compactness, or connectedness hypotheses. These conventional hypotheses clarify this possible reading rather than repair the source by fiat.
\subsection{Short English statement}
The fragment's printed ``Anosov automorphism'' has two readings. Hyperbolic toral automorphisms are mixing by the reported Fourier criterion; transitivity of arbitrary Anosov diffeomorphisms is the distinct general question reported as open. The fragment is not a complete new conjecture.
\subsection{Sources}
\begin{itemize}
\item American Institute of Mathematics, \emph{Emerging applications of measure rigidity}, fragment referring to Conjecture 3 and source bibliography labels [61, 14, 69, 90], \url{https://aimath.org/WWN/measrigid/measrigid.pdf}. The corresponding bibliography entries are not supplied.
\item The attached research aid identifies partial-result versions \texttt{arXiv:2410.15740v3}, revised 9 December 2025, and \texttt{arXiv:2410.12958v5}, revised 6 May 2026; no titles or authors are supplied here.
\item Local UnsolvedMath registry, permanent ID $20003266$, source code \texttt{catalogue entry}; canonical location \texttt{aim-other-notes.json}, \texttt{notes[72]}, dated 14 August 2026. Research attempt 2 is labelled context only and ``not a problem'' at run, with proved rigor and low novelty confidence; it is machine-generated, not peer reviewed, and requires independent verification. The registry's open label is retained as metadata, not a newly verified status.
\end{itemize}
\end{document}
